\section{Introduction}

Street-level imagery has become a standard instrument for measuring the built
environment, and the review literature already covers several hundred
applications \citep{biljecki2021street,ito2024understanding}. A growing share of
that work is longitudinal. Two images of one location, captured years apart, are
compared to establish what changed on the ground and what the change did to how
the street is perceived \citep{naik2017computer,stalder2024self,liu2025physical}.
Vision--language models have made the comparison cheap: one prompt returns a
change classification, a typology and a battery of perceptual scores, with no
task-specific training data. The results are usually published as maps in which
each sample point carries its own perception change, and spatial pattern is read
off those maps.

Figure~\ref{fig:teaser} is one such point. A vacant lot in Oakland was
photographed by the Street View car twelve times between 2008 and 2023. Eight
townhouses went up on it in 2017, and the instrument follows that without
difficulty: the street is judged wealthier and better maintained from the visit
at which the frame is finished, and it holds the new level for the six years
that follow. The same figure shows two other pairs of visits at that standpoint
in which nothing was built or removed. Between June 2021 and January 2023 the
same trees lose their leaves; between July and November 2015 the same vacant lot
is photographed under cloud and then in sun. Each of those moves the score by
three-quarters or more of what separates two different streets in the same city.
Nothing in the pipeline flags either reading as unusable. On a map, both would
appear as change.

\begin{figure}[!tb]
  \centering
  \includegraphics[width=\linewidth]{figures/fig01_teaser.pdf}
  \caption{One Oakland standpoint, twelve visits, fifteen years. (a) Three
  pairs of visits: the same vacant lot under cloud and then in sun, the lot
  rebuilt into eight townhouses, and the same trees in leaf and then bare. The
  frame is teal where the street changed and amber where only the photograph
  did. (b) What \texttt{gemini-2.5-flash} returned at each of the twelve
  visits, on three of the twelve dimensions; each bracket gives the mean
  absolute move over all twelve dimensions for the pair above it, and the
  shaded band spans the redevelopment that CityPulse marks. (c) The same three
  numbers placed against the distribution of that quantity for all 4{,}648
  control frame pairs, and for pairs of different standpoints in one city.}
  \label{fig:teaser}
\end{figure}

Whether that instance is representative is an empirical question, and it decides
whether the reporting scale now in use is supportable. We answer it directly, by
measuring what the instrument does when the place does not change. The panel for
that test is the set of Google Street View standpoints that the CityPulse
annotators \citep{huang2024citypulse} inspected epoch by epoch and marked as
never undergoing substantial redevelopment. Across 4{,}648 consecutive-epoch
frame pairs at 435 such standpoints, a second photograph moves a perception
score by two-thirds of what separates two different streets in the same city. A
single location's measured change is dominated by this quantity.

Two questions follow. The first is whether the movement is an acquisition bias
that could be modelled and subtracted. Six image statistics covering scattering,
contrast, colour, exposure, sharpness and specularity, entered as levels, as
differences and in a gradient-boosted model cross-fitted by standpoint, explain
almost none of the epoch-to-epoch variance, while the same six statistics
separate real weather conditions by up to two and a half corpus standard
deviations. The second is whether what remains is centred on zero, which decides
whether averaging helps. It very nearly is. The signed drift on control
standpoints reaches 0.12 points at most, a tenth of the between-street
difference; it is positive on almost every dimension; and it grows with the
interval between captures rather than with the calendar. That is the signature
of small physical change the labelling scheme does not record, and it sets a
floor on what aggregation can recover.

Acquisition conditions are not innocuous in general. A controlled experiment, in
which sixteen families of degradation are applied to fixed base images so the
scene is constant by construction, shows that degradation does move scores. It
also shows why the observational test came out flat: Google resamples every
panorama into one virtual camera before an analyst sees it, which removes the
resolution, projection and field-of-view variation that the controlled
experiment finds most consequential. Where the camera is free to vary, on
crowdsourced imagery, geometry alone makes a model report physical change in
45\% of pairs whose scene is identical by construction. The same experiment
shows that the sign of the acquisition response depends on which model is asked:
across three models scored on identical inputs, half the dimensions disagree on
sign, and one reverses from significantly negative under one model to
significantly positive under another.

Aggregation restores the measurement. Over a few hundred paired observations,
standpoints that cross a labelled change point are judged wealthier, better
maintained and more enclosed, and less green, and the result survives balancing
the comparison across the five cities. The variance
measured here converts directly into the number of paired observations an effect
of a given size requires, and that number is the operational result of the
paper.

We make four contributions.

\begin{enumerate}
\item A reliability ladder for vision--language perception scoring, placing
      sampling noise, encoding noise, prompt-order noise, same-place
      photographic noise and between-place variation on one scale
      (Section~\ref{sec:ladder}). Prompt structure turns out to be an instrument
      property of the same magnitude as image encoding, which bears on the
      twenty- and thirty-dimension batteries now in applied use.
\item A decomposition of the same-place residual into the part six acquisition
      statistics can predict, the part that is systematic, and the part that is
      neither, together with a measurement of how far those statistics can see
      (Section~\ref{sec:null}).
\item A three-model test of whether the sign of the acquisition response belongs
      to the task or to the model, on identical images under identical
      conditions (Section~\ref{sec:probe}), and a measurement of what camera
      geometry alone produces on crowdsourced imagery
      (Section~\ref{sec:optics}).
\item The sample size at which longitudinal perception measurement becomes
      usable, and the aggregate signal it recovers there
      (Section~\ref{sec:signal}).
\end{enumerate}

\section{Related work}

\paragraph{Longitudinal street-level measurement.}
\citet{naik2017computer} established the template: paired Google Street View
imagery from 2007 and 2014 across five US cities, a physical-change score, and a
test of urban theory against it. \citet{stalder2024self} learn embeddings
insensitive to lighting and season and sensitive to structural change, and
validate them on 1{,}449 manually labelled pairs before mapping London.
\citet{liu2025physical} combine a change-detection network with a perception
model trained on Place Pulse~2.0 to study New York City and Memphis from 2007 to
2023, and report hot-spot and cold-spot clusters of perception change at 300\,m
sample points. That is the closest published antecedent, and it reports at
exactly the scale examined here.

\paragraph{Acquisition conditions as a source of bias.}
That street-level imagery carries the conditions of its capture is established.
\citet{zhao2025seasonal} quantify seasonal bias across forty cities and show
that it changes a downstream clustering of urban function;
\citet{ki2026weather} use repeat Google Street View imagery to show that weather
at the moment of capture shifts perception scores, by different amounts in
different cities. Both work with perception models trained on human ratings. The
present paper asks the complementary question for a prompted vision--language
model, and separates the part of the same-place difference that whole-frame
acquisition summaries can predict from the part they cannot.
\citet{alpherts2025artifacts} and \citet{fan2025coverage} address a different
threat: the biases in where crowdsourced imagery exists at all and how densely
it samples the street.

\paragraph{Perception measurement and its reliability.}
The human ground truth most of this literature rests on comes from Place Pulse
\citep{salesses2013collaborative} and its successor Place Pulse~2.0
\citep{dubey2016deep}, and vision--language models are now used in their place.
\citet{mushkani2026benchmarks} evaluates seven such models
against 100 Montreal scenes annotated on thirty dimensions by twelve
participants, and argues that benchmarks should report inter-annotator
reliability alongside model alignment. That axis is agreement with human
raters. The axis measured here is the stability of a score under a second
photograph of the same scene, which is what a longitudinal design consumes and
which agreement with human raters does not bound.
\citet{he2026urbanfeel} benchmark twenty models on temporal change questions
over multi-temporal imagery from eleven cities, measuring answer accuracy.

\paragraph{Pairing and viewpoint.}
\citet{torkko2026pairwise} addresses the problem that metadata alone is
insufficient to find visually aligned pairs, and releases a tool that filters
candidates on feature matches and semantic-mask agreement. That filter operates
on viewpoint. Section~\ref{sec:optics} shows that once viewpoint is controlled,
camera geometry remains, and measures what it does on its own.

\section{Data and design}
\label{sec:data}

Four experiments appear in this paper. They differ in one respect that decides
how each result should be read: how much of the world is held still while the
measurement is repeated. Figure~\ref{fig:design} states that once, in order.

\begin{figure}[!tbp]
  \centering
  \includegraphics[width=\linewidth]{figures/fig02_design.pdf}
  \caption{The four experiments, ordered by how much they pin down. The lock
  marks what is clamped and the amber icons mark what is left free. The field
  panel holds one standpoint and lets a second visit change the light, the
  season and the traffic. The noise floor holds one photograph and varies the
  random seed, the JPEG quality and the order of the twelve questions. The
  perturbation experiment holds one scene and applies sixteen degradation
  families at 57 levels. The optics experiment holds one panorama and renders
  it through two virtual cameras, 61$^\circ$ against 96$^\circ$.}
  \label{fig:design}
\end{figure}

\paragraph{Units.}
Five terms recur and are not interchangeable. A \emph{panorama} is one Street
View capture, identified by a panorama identifier. A \emph{frame} is one image
rendered from a panorama at one heading; we request two headings per panorama. A
\emph{standpoint} is a kerbside position that CityPulse revisits across epochs,
holding a series of panoramas. A \emph{frame pair} joins two consecutive epochs
at one standpoint and one heading, and a \emph{transition} is the same step with
the two headings averaged. The reliability ladder is reported per frame pair,
because that is what an analyst who scores one image per location per epoch
obtains. Everything inferential is computed per transition, with standard errors
clustered by standpoint, because the two headings of one transition are repeated
measurements of the same step.

\paragraph{The longitudinal panel.}
CityPulse \citep{huang2024citypulse} publishes 9{,}751 Google Street View image
records, covering 8{,}846 distinct panorama identifiers at 757 standpoints in
five US cities from 2007 to 2023, each image inspected by hand and marked as a
change point or not. Requesting imagery by panorama identifier returns that
specific capture, so a pair contains no positional tolerance of the kind a
proximity match introduces. Re-driving is not perfectly repeatable: consecutive
captures at one standpoint sit a median of 1.4\,m apart, and 95\% of them within
5\,m, far tighter than the 10\,m gate used to match crowdsourced imagery in
Section~\ref{sec:optics}.

We re-checked every identifier against the Street View metadata endpoint in
August 2026. 51.2\% still resolve, and survival is scattered within series
rather than clustered, so 88.5\% of standpoints retain at least two epochs. The
endpoint also returns Google's own capture date, which disagrees with the
published index for 2.0\% of the surviving identifiers; since the date decides
which two images form a consecutive pair, we take the endpoint as authoritative
throughout. Each surviving panorama was requested twice through the Street View
Static API, at headings 0$^\circ$ and 180$^\circ$, 640 by 640 pixels, at the
default 90$^\circ$ field of view and zero pitch, so heading, field of view and
output size are fixed by the request and do not vary across epochs. That
returned 9{,}051 frames at 4{,}526 panoramas; 59 responses failed to parse, and
8{,}992 frames were scored.

Figure~\ref{fig:area} shows where the standpoints are and when they were
photographed. The control and labelled-change sets are drawn from the five
cities in very different proportions, which Section~\ref{sec:signal} takes into
account.

\begin{figure}[!tbp]
  \centering
  \includegraphics[width=\linewidth]{figures/fig03_study_area.pdf}
  \caption{The panel. (a) Standpoints in each city, separated into those that
  contribute control pairs and those that carry a labelled change point.
  (b) When each city was photographed; marker area is the number of scored
  frames in that year, and the right-hand figure is the median number of epochs
  per standpoint.}
  \label{fig:area}
\end{figure}

\paragraph{Corpus for the controlled experiments.}
Global Streetscapes \citep{hou2024global} provides 10.0 million street-level
images from 688 cities with per-image acquisition attributes, including a
manually annotated subset carrying human labels for weather, glare, reflection,
quality, platform, view direction and panorama status. Those labels serve twice:
to calibrate the perturbations against the conditions they are named for, and to
select clean base images.

\paragraph{Crowdsourced imagery.}
For the camera-geometry experiment we enumerate Mapillary imagery through the
vector-tile layer and retrieve per-image detail through the graph API. Four
tiles over San Francisco return 436{,}838 images, of which 28{,}160 grid cells at
20\,m by 45$^\circ$ resolution hold imagery from two or more distinct years.

\paragraph{What is scored.}
Twelve perception dimensions are used throughout: six aligned with Place
Pulse~2.0, for which human ground truth exists, and six physically auditable
dimensions a reader can check against the image. Each is anchored at 0, 5 and 10
and returned with a short evidence clause naming what in the frame drives the
score. Two dimensions, \emph{boring} and \emph{depressing}, are reverse-coded;
every figure flips them so that a positive shift always means the street looks
better. What the paper measures is the repeatability of these prompted scores,
which is a different question from their agreement with human raters.

The panel, the noise floor and the change-detection evaluation are all scored by
\texttt{gemini-2.5-flash} at temperature zero, so that one instrument runs
through the observational half of the paper. The controlled experiments of
Sections~\ref{sec:probe} and~\ref{sec:optics} vary the model deliberately and
name it wherever they do. The noise floor is measured on the 150 base images
that also serve the perturbation experiment, since it is a property of the
instrument rather than of the panel.

Throughout the paper, differences in score are reported against one yardstick:
the mean absolute difference between two randomly chosen different standpoints
\emph{in the same city}, 1.19 points on the ten-point scale. That is the
variation a map of one city exists to display. Where a figure reports an effect
in standard-deviation units instead, the standard deviation is that of the same
standpoint means.

\section{A second photograph of the same street}
\label{sec:ladder}

Five conditions bracket what a longitudinal comparison actually contains. Ask
the same question of the same image twice, and only the sampling of the model's
own output changes. Re-encode that image at a slightly different JPEG quality,
and the pixels move without the scene moving. Reorder the twelve questions in
the prompt and leave the image alone, and only the instrument changes.
Photograph the same unchanged street on a second visit, and everything a second
visit brings changes with it. Compare two different streets, and the scene
changes as well. Figure~\ref{fig:ladder} puts all five on one scale. Three
readings follow.

\begin{figure}[!tbp]
  \centering
  \includegraphics[width=\linewidth]{figures/fig04_ladder.pdf}
  \caption{(a) Mean absolute difference in perception score under four repeat
  conditions, averaged over twelve dimensions. Each sits in a track as wide as
  the difference between two standpoints in the same city, 1.19 points, so the
  filled share is the quantity in the right-hand column. (b) The same
  comparison dimension by dimension; the percentage is the same-place
  difference as a share of that dimension's own between-place difference.}
  \label{fig:ladder}
\end{figure}

The model is close to deterministic at temperature zero. Repeated calls on one
image differ by less than one part in a hundred of the between-place difference,
so averaging repeated samples buys nothing, and any stability obtained that way
is illusory. Sampling variance explains none of what follows.

Prompt structure is an instrument property of the same order as image encoding.
Re-encoding an image at a JPEG quality four points away, which produces a
visually indistinguishable file, moves a score by a fifth of the between-place
difference. Reordering the twelve questions while leaving the image untouched
moves it by the same amount, with \emph{safe} the most position-sensitive
dimension of the twelve. Applied work now routinely asks for twenty or thirty
dimensions in one pass; a battery of that length is a longer list to permute,
and the effect is not currently reported.

The dominant term is none of these. A second photograph of the same unchanged
standpoint differs by three times the prompt-order effect and two-thirds of the
between-place difference, and Figure~\ref{fig:ladder}b shows this is not an
artefact of averaging over dimensions. Eleven of the twelve sit between 41\% and
89\% of their own between-place difference, and the twelfth, \emph{lively},
crosses it: two photographs of the same street differ on how lively it looks by
more than two different streets in the same city do. The variation is spread across the
panel rather than concentrated in a few unstable standpoints, with only 7\% to
14\% of the variance in the pairwise difference lying between standpoints.
Averaging the two headings of a standpoint, which costs one extra request per
epoch, brings the same-place difference down by a fifth, to 0.64 points.

\subsection{What the residual is made of}
\label{sec:null}

For every image we compute six statistics: the dark-channel mean
\citep{he2011single}, whole-frame RMS contrast, mean saturation, mean luminance,
mean absolute Laplacian response, and the share of near-saturated pixels. These
capture scattering, contrast, colour, exposure, sharpness and specularity, and
they are the statistics with which the perturbations of
Section~\ref{sec:probe} are calibrated.

Figure~\ref{fig:null}a fits them to the epoch-to-epoch difference on the control
standpoints. Across all six statistics and all twelve dimensions the largest
single correlation is 0.19, and a gradient-boosted model over the levels and
differences of all six, plus the interval between captures and the seasonal
offset, cross-fitted with folds assigned by standpoint, returns a negative
out-of-fold $R^2$ on ten of twelve dimensions. \emph{Greenery} is the exception,
at 0.075. There is no acquisition component of consequence to subtract.

\begin{figure}[!tbp]
  \centering
  \includegraphics[width=\linewidth]{figures/fig05_null.pdf}
  \caption{(a) Out-of-fold $R^2$ of a gradient-boosted model over six
  acquisition statistics, fitted to the epoch-to-epoch difference on 4{,}648
  control frame pairs, with folds assigned by standpoint. (b) The same six
  statistics on 4{,}000 manually annotated Global Streetscapes images, grouped
  by human label and expressed as a displacement from the clear-weather mean in
  corpus standard deviations.}
  \label{fig:null}
\end{figure}

A null is worth only as much as the covariates behind it can see, and
Figure~\ref{fig:null}b measures that on 4{,}000 manually annotated images. Fog
displaces the dark channel by 2.4 corpus standard deviations, cuts saturation to
a third and multiplies the specular fraction eightfold; cloud, snow and rain
each sit in a distinct position; images annotated for glare are displaced by 1.2
standard deviations. Windscreen reflection is the exception. Images annotated as
carrying a reflection have the same profile as those annotated without one on
all six statistics, because reflection is a localised, spatially structured
artefact and a whole-frame mean averages it away. Reflection is also the
acquisition attribute whose prevalence grew fastest over the study period. The
imagery used here is rendered by Google from a fixed virtual camera, which
removes the geometric part of that artefact but not a reflection present in the
source panorama.

What is left has a small systematic part. Measured on the 2{,}367 control
transitions with standard errors clustered by standpoint, the signed drift
reaches 0.118 points on \emph{wealthy} and 0.110 on \emph{building condition},
under a tenth of the between-place difference in both cases, and it is
distinguishable from zero on eleven of twelve dimensions. Its sign is the same
throughout: a later photograph of the same unchanged standpoint is judged
marginally safer, wealthier, better maintained and more attractive.

The timing says what produces it. The drift grows with the interval between the
two captures ($r = 0.088$, $p < 10^{-4}$) and is flat against the calendar
($r = -0.008$, $p = 0.70$), so it is not a secular improvement in Street View
imagery. It is what a street does over time when nobody records it: the mean
absolute difference in the same data rises from 0.75 points under a year to 0.91
points at four to eight years, and Section~\ref{sec:threshold} finds real
physical change among control pairs when they are inspected one by one. The
practical consequence is a floor. Averaging removes the random part of the
same-place difference and leaves about a tenth of a point of drift, which is
below the effects Section~\ref{sec:signal} detects and comparable to the ones it
does not.

\subsection{What that does to reporting scale}

An estimate at a single location carries a standard error of the same order as
the between-place standard deviation. A map in which each sample point carries
its own perception change is therefore displaying, at the point level, a
quantity whose noise is comparable to the cross-sectional signal it is being
compared against. Aggregation is the condition under which such a design
produces anything, and Section~\ref{sec:signal} gives the sample size at which
it does.

\section{What acquisition does when it is free to vary}
\label{sec:probe}

The flat result in Section~\ref{sec:null} is a statement about rendered Street
View imagery. The complementary question is what acquisition conditions do when
they are allowed to move, and that requires holding the scene fixed by
construction rather than by a label.

We apply sixteen families of perturbation at 57 levels to 150 base images
restricted to daylight, clear-weather, good-quality, non-panoramic frames and
stratified across six continents. Each family reproduces one axis along which
the corpus varies. Atmospheric scattering follows the Koschmieder model with the
dark-channel prior \citep{he2011single} as its calibration statistic, and camera
yaw is applied as the exact homography $K R K^{-1}$ for a pinhole source and by
resampling the sphere for an equirectangular one, so the manipulation is exact
in both cases rather than approximated by an image shift. Perturbation strengths
are set so that each reproduces the statistical signature of the corresponding
real condition, measured on the annotated corpus of Figure~\ref{fig:null}b. The
fitted fog strength is $\beta = 0.95$; the $\beta = 2.0$ level is roughly twice
real fog and is reported as an upper bound.

\begin{figure}[!tbp]
  \centering
  \includegraphics[width=\linewidth]{figures/fig06_surface.pdf}
  \caption{(a) Perturbation family against perception dimension for
  \texttt{qwen2.5-vl-72b-instruct}, at the level where each family produces its
  largest mean absolute response, in units of the between-place standard
  deviation, with that magnitude in the gutter. Because the level is chosen on
  the response, each row is an upper bound on its family. (b) The
  dose--response for atmospheric scattering, with bootstrap intervals and the
  calibrated strength of real fog marked.}
  \label{fig:surface}
\end{figure}

\subsection{Magnitude}

Figure~\ref{fig:surface}a gives the response surface for
\texttt{qwen2.5-vl-72b-instruct}, one of the models used by published
urban-perception benchmarks \citep{mushkani2026benchmarks}. Each row is shown at
the level where that family's mean absolute response is largest, so a row that
stays inside the noise floor stays inside it at its own most damaging setting. A
cell has to pass two tests to count: a bootstrap interval that excludes zero,
and a shift larger than the noise floor that prompt order and re-encoding set
for that dimension in Figure~\ref{fig:ladder}. Taken one cell at a time the
picture is restrictive: 29 of 192 clear both. Atmospheric scattering dominates,
followed by motion blur, defocus, exposure and camera yaw. Ten families clear
nothing anywhere: resolution, JPEG quality, white balance, saturation, contrast,
glare, reflection, foliage state, field of view and rain.

Dose--response relationships are far cleaner than single cells, because model
variability is not ordered by perturbation strength while the response is.
Figure~\ref{fig:surface}b shows the largest family: at the calibrated strength
of real fog the shift is about a fifth of a between-place standard deviation on
the most affected dimensions, and it grows monotonically with dose.

\subsection{Direction, and whose property it is}

Effect size and direction answer different questions. A family that shifts a
dimension by less than the noise floor still carries information about the sign,
and sixteen families give sixteen readings of it. We take the sign in each
family for three models scored on the identical 150 base images, restricted to
the 34 perturbation levels that all three runs share, so that the model is the
only thing that differs.

\begin{figure}[!tbp]
  \centering
  \includegraphics[width=\linewidth]{figures/fig07_models.pdf}
  \caption{(a) Mean shift over sixteen degradation families, by dimension and by
  model, on identical images under identical conditions. Filled markers are
  significant at 5\% on a two-sided binomial test over the sixteen family signs.
  (b) Correlation between the three models' full family-by-dimension response.
  None of these three is the \texttt{gemini-2.5-flash} instrument used for the
  field panel.}
  \label{fig:models}
\end{figure}

In Figure~\ref{fig:models}a, six of twelve dimensions have all three models
agreeing on the sign of the response, and no dimension is significant in all
three. One disagreement reverses a significant result: \emph{beautiful} is
significantly negative under \texttt{gemini-3.7-flash} and significantly
positive under \texttt{llama-4-maverick}. \emph{Boring} reverses in sign as
well, with one of the two models individually significant. A study reporting
\emph{beautiful}, with either model, would have had a result at conventional
significance whose sign the choice of model determined. What agreement exists is
concentrated in one pair (Figure~\ref{fig:models}b): the qwen--gemini
correlation over the full family-by-dimension response is 0.59, against 0.42 and
0.28 for the pairs involving llama.

The sign of the acquisition response therefore depends on which model is asked.
The consequence is practical: a directional correction cannot be applied without
first fixing the model, and a correction estimated with one model may point the
other way under another.

\section{Camera geometry in crowdsourced imagery}
\label{sec:optics}

Crowdsourced imagery has no panorama identifier to request, so a pair must be
identified from metadata. Mapillary exposes structure-from-motion-refined
position and heading, and a field named \texttt{merge\_cc}, documented as the
connected component of images aligned into a common reconstruction. The refined
fields are necessary and available: present on every image checked, moving the
raw position by a median of 7.1\,m and the raw compass heading by a median of
3.1$^\circ$. The connected-component field does not do what its name suggests.
Only 15.0\% of cross-year pairs share one, the share is unrelated to how far
apart the two images are, and 40 frames from a single sequence captured seconds
apart by one moving camera carry 22 to 28 distinct values. It cannot serve as a
co-location key.

Applying the geometric criterion alone, 47.7\% of cross-year candidate pairs fall
within 10\,m and 20$^\circ$. Visual inspection of the pairs that pass shows the
remaining obstacle is the camera. A 2015 dashcam covers 61$^\circ$ by
48$^\circ$, a 2024 action camera 96$^\circ$ by 64$^\circ$, and an
equirectangular panorama the whole sphere. Two frames from the same point facing
the same way frame nothing alike, and a panorama has no heading in the sense the
criterion assumes.

We therefore resample both frames into one virtual pinhole camera, with a shared
field of view derived from what the two sources can each cover and a yaw equal to
the difference between the target heading and each image's own refined heading.
Working from the relative heading avoids depending on the orientation convention
of the reconstruction frame. Deriving the field of view per pair rather than
fixing it matters: a fixed 80$^\circ$ target leaves 64\% of the output frame
outside the source, while the derived target leaves under 1\%.

\begin{figure}[!tbp]
  \centering
  \includegraphics[width=\linewidth]{figures/fig08_optics.pdf}
  \caption{(a) One Mapillary panorama rendered as two virtual cameras, above,
  and normalised into a common camera, below; the wedge gives the two fields of
  view in plan. (b) Over 40 such panoramas, the share on which the model
  reported a physical change, with the mean absolute perception difference
  alongside. The two cameras differ by 61$\times$48$^\circ$ against
  96$\times$64$^\circ$ and 4$^\circ$ of yaw; the normalised row keeps
  3.1$^\circ$ of yaw and 6\% of field as residual error, and the last row
  renders both frames with identical parameters. Both frames of every pair come
  from a single panorama, so the scene is identical by construction.}
  \label{fig:optics}
\end{figure}

Figure~\ref{fig:optics} isolates the effect. One panorama is rendered twice,
once as the narrow camera and once as the wide camera offset by 4$^\circ$. Both
frames show the same scene, from the same point, at the same instant, under the
same light, so any reported change is an artefact. Under raw optics the model
reports physical change on 45\% of the forty panoramas, naming building facade
change six times, amenities and furniture four, road and traffic infrastructure
four, storefront signage three and boundary configuration three, on scenes that
did not change. Normalising both frames into one virtual camera, while keeping
the residual alignment error that survives in the field, reduces that sixfold.

The two outcomes separate. Change classification is largely fixed by
normalisation; the perception artefact is not, falling only from 0.39 to 0.35,
because perception scores remain sensitive to three degrees of yaw and six
percent of field of view, which is below what any alignment can guarantee.
Feature matching is not improved at all: median inlier counts are 14 before and
13 after, with only 30\% of pairs improving, because scale-invariant keypoint
matching is already invariant to the differences normalisation removes and the
resampling costs detail. Keypoint matching and a vision--language model ask
different questions of the same pair, and a pipeline validated on the first is
not thereby validated on the second.

\section{What survives aggregation}
\label{sec:signal}

Averaging removes the random part of the same-place difference and leaves the
drift of Section~\ref{sec:null}, so the question is how many paired observations
are needed before an effect larger than that drift is resolvable.
Figure~\ref{fig:signal}a compares the 224 transitions spanning a labelled change
point, at 193 standpoints, against the 2{,}367 control transitions at 435
standpoints.

\begin{figure}[!tbp]
  \centering
  \includegraphics[width=\linewidth]{figures/fig09_signal.pdf}
  \caption{(a) Perception change across a labelled change point against control
  transitions, in score points, with 95\% intervals clustered by standpoint.
  Filled circles are the pooled contrast; open diamonds re-weight the five
  cities so that the two groups have the same city composition. (b) The
  labelled-change transitions required to detect an effect of a given size at
  80\% power, holding the control set at its realised size.}
  \label{fig:signal}
\end{figure}

Four of the twelve dimensions are distinguishable from the control transitions,
and they describe a coherent process: a standpoint that crosses a labelled
change point is judged wealthier, better maintained and more enclosed, and less
green. Redevelopment adds built fabric and removes vegetation, and the
instrument reports both. Three of the four are physically auditable dimensions
and one is Place-Pulse aligned, so the result does not rest on the perceptual
half of the battery alone.

That comparison is exposed to a compositional difference. Oakland supplies 202
control standpoints and 19 labelled-change ones, Los Angeles 54 and 104, so a
pooled contrast partly compares cities. Re-weighting the five cities so that the
two groups match, all four dimensions survive: greenery falls by 0.35, building
condition rises by 0.24, wealthy by 0.20 and enclosure by 0.19, each at
$p < 0.01$. The effects shrink by a fifth to a third and the pattern is the
same. Greenery is the one dimension that behaves differently in one place,
falling in Boston, Los Angeles and San Francisco and rising in Seattle.

Figure~\ref{fig:signal}b converts the measured variance into sample size. An
effect the size of the shifts that clear in panel (a) needs about 88
labelled-change transitions against a control set of the size realised here, and
224 are available. An effect of a tenth of a point needs 1{,}010, and one of six
hundredths needs 6{,}240. Below about a tenth of a point the drift of
Section~\ref{sec:null} becomes the binding constraint rather than the sample
size, because a systematic tenth of a point does not average away. A study
reporting perception change at a single location is working three orders of
magnitude below the smaller of these two thresholds.

\subsection{Where the change threshold sits}
\label{sec:threshold}

Part of the residual in Section~\ref{sec:null} is real change that no label
records, and the same ambiguity governs how a change rate should be read.
Table~\ref{tab:threshold} evaluates paired change detection against the
CityPulse labels. The pairs are the consecutive-epoch pairs at heading
0$^\circ$ for the standpoints whose series carries a labelled change point, one
pair per step, giving 172 pairs of which 24 are labelled change points. The
three prompts differ only in what they specify; the images and the model are
identical throughout.

\begin{table}[!tbp]
  \caption{Paired change detection against the CityPulse labels, on 172
  consecutive-epoch pairs. The reference label marks substantial redevelopment;
  the prompts differ only in what they specify.}
  \label{tab:threshold}
  \centering
  \small
  \begin{tabular}{lrrrrrrr}
    \toprule
    Prompt & Reports change & TP & FP & TN & Precision & Recall & Specificity \\
    \midrule
    Task description only                & 97.7\% & 24 & 144 & 4  & 0.143 & 1.000 & 0.027 \\
    + acquisition artefacts to disregard & 86.0\% & 23 & 125 & 23 & 0.155 & 0.958 & 0.155 \\
    + threshold for substantial change   & 76.7\% & 22 & 110 & 38 & 0.167 & 0.917 & 0.257 \\
    \bottomrule
  \end{tabular}
\end{table}

\begin{figure}[!tbp]
  \centering
  \includegraphics[width=\linewidth]{figures/fig10_threshold.pdf}
  \caption{Three of the six false positives with the smallest acquisition
  difference and the shortest interval. CityPulse marks none of these
  standpoints as a change point over the whole series, and each pair carries a
  persistent physical change: a construction hoarding erected along the fence
  line, fresh graffiti with repainted crossing markings, and a hoarding with a
  newly painted red kerb.}
  \label{fig:threshold}
\end{figure}

Specificity rises ninefold across the three prompts while recall falls only from
1.000 to 0.917. Precision stays near 0.16 throughout, because the model keeps
reporting change on three-quarters of the pairs while the reference label marks
it on 14\%. Some of that gap is a difference in definition rather than model
error. Figure~\ref{fig:threshold} shows three of the six pairs with the
smallest acquisition difference and the shortest interval; each carries a
physical change that persists and that the reference label does not record.
CityPulse labelled substantial redevelopment;
the model was asked about any persistent physical change; over a four-month
interval a street does change in small ways. Across the three prompts the model
calls change on between 76.7\% and 97.7\% of these pairs while the reference
prevalence is 14.0\%, and the operational threshold that separates the two is a
free parameter this literature does not state.

\section{Discussion}

\paragraph{Reporting scale.}
The result that governs the rest is that a second photograph of an unchanged
street reproduces two-thirds of the difference between two different streets in
the same city. Nothing in the pipeline removes this: neither sampling variance,
nor prompt order, nor the acquisition statistics that separate real weather
conditions cleanly. Point-level maps of perception change are displaying it. The
usable unit is a few hundred paired observations, and the field currently
reports at a unit of one.

\paragraph{Where correction is and is not needed.}
The absence of a predictable acquisition component in the Street View panel is a
property of that imagery, not a general finding. Google renders every panorama
into a fixed virtual camera before an analyst sees it, and
Section~\ref{sec:optics} shows what happens without that step. On crowdsourced
imagery, rendering both frames into a common camera removes most of the spurious
change classification and little of the perception artefact, so it is worth
doing and it is not sufficient on its own.

\paragraph{Model choice.}
Across three models on identical images, six of twelve dimensions agree on the
sign of the acquisition response and none is significant in all three.
\emph{Beautiful} reverses from significantly negative under one model to
significantly positive under another. The sign belongs to the model, which means
a directional correction cannot be applied without fixing the model first.

\paragraph{Reproducibility of the imagery.}
Half the panorama identifiers published with a January 2024 dataset no longer
resolved in August 2026, and Google's own capture date disagreed with the
published one for a further 2.0\% of those that did. A study identifying its
imagery by Street View panorama identifier cannot retrieve the images it used.
The release accompanying this paper therefore carries identifiers, retrieval
dates, the capture date the endpoint returned, the prompts and the derived
scores, which is what remains distributable.

\section{Limitations}

The reliability measurement is made on Google Street View imagery, rendered to a
common virtual camera and requested at a fixed heading and field of view, so the
same-place difference reported here is a lower bound for crowdsourced imagery,
where camera heterogeneity adds the component measured in
Section~\ref{sec:optics}. Control standpoints are certified by human labels that
mark substantial redevelopment, and Section~\ref{sec:threshold} shows that
streets change in ways such a scheme does not record; the signed drift in
Section~\ref{sec:null} is the size of that contribution, and separating it
completely from instrument noise would require a labelling scheme with an
explicit magnitude threshold applied to the control set. The panel covers five
US cities, and the three models of Section~\ref{sec:probe} are not the model
that scored the field panel, so the sign disagreement shows that model choice
matters, and does not supply a correction table for \texttt{gemini-2.5-flash}. Two of the sixteen perturbation families cannot be
calibrated against what their names refer to, since images annotated for glare
carry the signature of low-sun illumination rather than of a source in frame,
and images annotated for reflection have the same whole-frame profile as those
without one. Base images reach 3.15 megapixels at the median, because the Global
Streetscapes distribution carries 2048-pixel thumbnails, so the finding that
resolution produces no effect above the noise floor is established over 0.5 to
3.2 megapixels rather than over the full observed range.
